\begin{appendixtable}
\centering
\small
\caption{Qualitative coding themes.}
\label{tab:qualitative-codebook}
\begin{singlespace}
\renewcommand{\arraystretch}{1.08}
\begin{tabular}{@{}>{\raggedright\arraybackslash}p{1.55in}>{\raggedright\arraybackslash}p{4.25in}@{}}
\hline
\textbf{Theme} & \textbf{Definition} \\
\hline
Semantic fit / meaning clarity & The response connects the vibration to the meaning of the visual event, such as success, error, notification, completion, or action outcome. \\
Timing and rhythm & The response discusses temporal structure, rhythm, pulse pattern, onset timing, repetition, delay, or synchronization with the visual sequence. \\
Intensity / salience & The response discusses vibration strength, force, sharpness, attention capture, or whether the stimulus is noticeable enough. \\
Smoothness / comfort & The response discusses pleasantness, comfort, smoothness, harshness, irritation, or bodily feel. \\
Randomness / lack of structure & The response describes a stimulus as random, noisy, chaotic, unstructured, arbitrary, or hard to interpret. \\
Progress / process feeling & The response links vibration to an ongoing process, loading, pulling, scrolling, transition, build-up, release, or completion over time. \\
Subtlety / appropriateness & The response discusses whether the vibration is appropriately restrained or excessive for the context. \\
General preference / dislike & The response expresses liking, disliking, preference, or ranking without a more specific reason. \\
\hline
\end{tabular}
\end{singlespace}
\end{appendixtable}
